\SourceLeaf{231}%
\Article{6}
\indent Sicuti, per translatam directionem normalis in superficiem curvam ad superficiem sphaerae, cuius puncto determinato prioris superficiei respondet punctum determinatum in posteriore, ita etiam quaevis linea, vel quaevis figura in illa repraesentabitur per lineam vel figuram correspondentem in hac. In comparatione duarum figurarum hoc modo sibi mutuo correspondentium, quarum altera quasi imago alterius erit, duo momenta sunt respicienda, alterum, quatenus sola quantitas consideratur, alterum, quatenus abstrahendo a relationibus quantitativis solum situm contemplamur.

Momentum primum basis erit quarundam notionum, quas in doctrinam de superficiebus curvis recipere utile videtur. Scilicet cuilibet parti superficiei curvae limitibus determinatis cinctae \textit{curvaturam totalem} seu \textit{integram} adscribemus, quae per aream figurae illi in superficie sphaerica respondentem exprimetur. Ab hac curvatura integra probe distinguenda est curvatura quasi specifica, quam nos \textit{mensuram curvaturae} vocabimus: haec posterior ad \textit{punctum} superficiei refertur, et denotabit quotientem qui oritur, dum curvatura integra elementi superficialis puncto adiacentis per aream ipsius elementi dividitur, et proin indicat rationem arearum infinite parvarum in superficie curva et in superficie sphaerica sibi mutuo respondentium. Utilitas harum innovationum per ea, quae in posterum a nobis explicabuntur, abunde, ut speramus, sancietur. Quod vero attinet ad terminologiam, imprimis prospiciendum esse duximus, ut omnis ambiguitas arceatur, quapropter haud congruum putavimus, analogiam terminologiae in doctrina de lineis curvis planis vulgo receptam (etsi non omnibus probatam) stricte sequi, secundum quam mensura curvaturae simpliciter audire debuisset curvatura, curvatura integra autem amplitudo. Sed quidni in verbis faciles esse liceret, dummodo res non sint inanes, neque dictio interpretationi erroneae obnoxia?

Situs figurae in superficie sphaerica vel similis esse potest situi figurae respondentis in superficie curva, vel oppositus (inversus); casus prior locum habet, ubi binae lineae in superficie curva ab eodem puncto directionibus inaequalibus sed non oppositis proficiscentes repraesentantur in superficie sphaerica per lineas similiter iacentes, puta ubi imago lineae ad dextram iacentis ipsa est ad dextram; casus posterior, ubi contrarium valet. Hos duos casus per \textit{signum} mensurae curvaturae vel positivum vel negativum distinguemus. Sed manifesto haec distinctio eatenus tantum locum habere potest, quatenus in utraque superficie pla%
